\documentclass[11pt,a4paper]{article}

\usepackage[margin=2.35cm,headheight=15pt]{geometry}
\usepackage{fontspec}
\usepackage{xeCJK}
\setCJKmainfont{Noto Serif CJK SC}[AutoFakeSlant=true]
\setCJKsansfont{Noto Sans CJK SC}[AutoFakeSlant=true]
\setCJKmonofont{Noto Sans CJK SC}
\xeCJKsetup{PunctStyle=kaiming}

\usepackage{amsmath,amssymb,amsthm,mathtools,bm}
\usepackage{booktabs}
\usepackage{enumitem}
\usepackage{xcolor}
\usepackage{tikz}
\usetikzlibrary{arrows.meta,positioning,calc,decorations.pathreplacing}
\usepackage[most]{tcolorbox}
\usepackage{hyperref}
\usepackage{fancyhdr}
\usepackage{titlesec}
\usepackage{array}
\usepackage{microtype}

\hypersetup{
  colorlinks=true,
  linkcolor=blue!50!black,
  urlcolor=blue!60!black,
  citecolor=blue!50!black
}

\definecolor{scblue}{RGB}{10,90,90}
\definecolor{scsoft}{RGB}{230,246,246}
\definecolor{warnbg}{RGB}{255,246,230}
\definecolor{okbg}{RGB}{232,246,236}

\theoremstyle{definition}
\newtheorem{definition}{定义}[section]
\newtheorem{example}{例子}[section]
\theoremstyle{plain}
\newtheorem{theorem}{定理}[section]
\newtheorem{proposition}{命题}[section]
\theoremstyle{remark}
\newtheorem{remark}{注记}[section]

\newtcolorbox{keybox}[1][]{
  colback=scsoft,colframe=scblue,fonttitle=\bfseries,
  title={核心要点},#1
}
\newtcolorbox{warnbox}[1][]{
  colback=warnbg,colframe=orange!70!black,fonttitle=\bfseries,
  title={容易混淆},#1
}
\newtcolorbox{hearbox}[1][]{
  colback=okbg,colframe=green!40!black,fonttitle=\bfseries,
  title={听课提示},#1
}

\pagestyle{fancy}
\fancyhf{}
\lhead{复旦粒子物理暑期学校 2026}
\rhead{SCET / Matthias Neubert}
\cfoot{\thepage}
\titleformat{\section}{\Large\bfseries\color{scblue}}{\thesection}{0.8em}{}
\titleformat{\subsection}{\large\bfseries\color{scblue!80!black}}{\thesubsection}{0.7em}{}

\usepackage{slashed}
\newcommand{\nbar}{\bar n}
\newcommand{\slashn}{\slashed{n}}
\newcommand{\slashnb}{\slashed{\bar n}}
\newcommand{\dd}{\mathrm{d}}
\newcommand{\order}{\mathcal{O}}
\newcommand{\scet}{\mathrm{SCET}}

\begin{document}

\begin{center}
{\LARGE\bfseries 软共线有效理论课前讲义}\\[0.45em]
{\large 复旦粒子物理暑期学校 2026}\\[0.25em]
{\large SCET / Soft-Collinear Effective Theory \quad Matthias Neubert (Mainz)}\\[0.8em]
{\normalsize 课前预习稿　从运动学与 EFT 写起　推导写全　英文术语单列}\\[0.35em]
{\small 依据 Indico 课表：8 月 17--21 日共五讲。本讲义非正式课堂讲义。}
\end{center}

\tableofcontents
\newpage

%======================================================================
\section{如何使用本讲义}
%======================================================================

SCET 解决的问题可以一句话概括：

\begin{quote}
QCD 过程里同时出现几个相距很远的能量尺度时，普通微扰论会被大对数毁掉。SCET 把\textbf{硬、共线、软}自由度写成不同场，用幂次计数做展开，再用重整化群把大对数求和回来。
\end{quote}

Matthias Neubert 是这一语言的主要建设者之一。他的课通常从 Sudakov 问题与有效场论讲起，而不是一上来写完整拉氏量。本讲义按他常见的逻辑铺路：散射运动学 $\to$ 尺度分离 $\to$ EFT $\to$ 软/共线极限 $\to$ 光锥坐标 $\to$ 幂次计数 $\to$ SCET 的基本结构 $\to$ 因子化与重求和。

五讲课表（Asia/Shanghai）：

\begin{center}
\begin{tabular}{lll}
\toprule
日期 & 时间 & 内容 \\
\midrule
8 月 17 日（一） & 14:50--15:50 & SCET \\
8 月 18 日（二） & 14:50--15:50 & SCET \\
8 月 19 日（三） & 16:10--17:10 & SCET \\
8 月 20 日（四） & 14:50--15:50 & SCET \\
8 月 21 日（五） & 14:50--15:50 & SCET \\
\bottomrule
\end{tabular}
\end{center}

\begin{hearbox}
Neubert 的板书会很快进入 $n^\mu,\bar n^\mu$ 与 $\lambda$ 计数。课前务必自己算一遍：一个沿 $z$ 轴超相对论运动的粒子，其 $p^+=n\cdot p$ 与 $p^-=\bar n\cdot p$ 哪个大、哪个小。把这件小事搞反，后面所有幂次都会错。
\end{hearbox}

%======================================================================
\section{预备知识：QFT 与散射运动学}
\label{sec:prereq}
%======================================================================

\subsection{需要已经会的}

\begin{itemize}[leftmargin=1.4em]
  \item 费曼规则、重整化、$\overline{\mathrm{MS}}$ 跑动耦合 $\alpha_s(\mu)$。
  \item 树图与一圈图的数量级：每个圈带 $\alpha_s$，还可能带对数。
  \item Dirac 旋量：$\gamma$ 矩阵、Dirac 方程 $(\slashed{p}-m)u=0$。质量可以先当作比硬标度小的量。
\end{itemize}

本讲义在欧氏/闵氏之间切换时会说明；QCD 振幅默认闵氏度规 $(+,-,-,-)$。

\subsection{两体与喷射运动学}

考虑 $e^+e^-\to$ 强子，质心能量 $Q$。一个两喷射事件里，每个喷射的不变质量 $M_J$ 满足
\begin{equation}
  M_J^2\ll Q^2.
\end{equation}
若再要求喷射很窄，内部还有更软的辐射。于是出现至少三个尺度：
\begin{equation}
\label{eq:scales}
  Q\ \gg\ M_J\ \gg\ \Lambda_{\mathrm{QCD}}
  \qquad\text{或更细地}\qquad
  Q\ \gg\ p_\perp\ \gg\ p_\perp^2/Q.
\end{equation}
末式正是 Sudakov 问题的标准层次：硬标度、共线标度、软标度。

对撞机过程 $pp\to Z+j$ 等同样如此：部分子动量分数、$p_T$、喷射半径 $R$、软辐射能量可以差好几个数量级。

\subsection{为什么``只算到固定阶''会失败}

典型截面在微扰论中长成
\begin{equation}
\label{eq:sudakov-series}
  \sigma\sim \sigma_0\sum_{n}\sum_{m\le 2n}
  c_{nm}\,\alpha_s^n\log^m\Bigl(\frac{Q^2}{M^2}\Bigr).
\end{equation}
对数来自软、共线相空间的积分。当 $\alpha_s\log^2(Q^2/M^2)\sim 1$ 时，后面的项并不比前面小，固定阶展开崩溃。必须对大对数做\textbf{重求和}（resummation）。SCET 的工程目标之一，就是把这件事变成解 RG 方程。

\begin{example}[Sudakov 形状因子]
一个带电粒子在硬动量转移 $Q$ 下几乎不辐射地保持其方向，振幅被指数抑制：
\begin{equation}
  F(Q^2,\mu^2)\sim \exp\Bigl[-\frac{\alpha_s C_F}{4\pi}\log^2\frac{Q^2}{\mu^2}+\cdots\Bigr].
\end{equation}
双对数 $\log^2$ 是软与共线叠加的指纹。Neubert 几乎一定从类似对象切入。
\end{example}

%======================================================================
\section{有效场论的基本思想}
\label{sec:eft}
%======================================================================

\subsection{积分掉重模式}

设理论有重标度 $M$ 与轻标度 $m$，且我们只问能量 $\sim m$ 的过程。把重场（或硬圈动量）积掉，得到
\begin{equation}
  \mathcal{L}_{\mathrm{EFT}}=\mathcal{L}_{\mathrm{light}}+\sum_i \frac{C_i(\mu)}{M^{\Delta_i-d}}\,\mathcal{O}_i,
\end{equation}
其中 Wilson 系数 $C_i$ 吸收短距物理，$\mathcal{O}_i$ 由轻场组成。

匹配（matching）：在能同时用全理论与 EFT 算的尺度上，让物理量相等，定 $C_i$。其后用 EFT 的 RG 把 $C_i$ 演化到低标度。

\subsection{EFT 能做什么}

\begin{enumerate}[leftmargin=1.6em]
  \item \textbf{幂次展开}：按 $m/M$ 系统展开，误差可估计。
  \item \textbf{尺度分离}：每个函数只依赖一个标度，对数不再混在一起。
  \item \textbf{对称性}：EFT 拉氏量必须实现全理论在低能仍存在的对称（规范、手征、重夸克自旋对称等）。
\end{enumerate}

HQET、ChPT、Fermi 理论都是这一套路。SCET 的特殊处在于：低能自由度不是``一种轻粒子''，而是\textbf{按方向与虚度分类的多种模式}。

\begin{keybox}
普通 EFT：重 vs.\ 轻。\\
SCET：硬 vs.\ 沿某个方向的共线 vs.\ 软（以及超软）。\\
同一个夸克场，在全 QCD 里只有一个；在 SCET 里被拆成 $\xi_n+\xi_{\bar n}+q_{us}+\cdots$。
\end{keybox}

%======================================================================
\section{软极限与共线极限}
\label{sec:limits}
%======================================================================

\subsection{共线}

一个质量为 $0$ 的部分子动量 $p^\mu$，若它几乎平行于某个类光方向 $n^\mu$，则
\begin{equation}
  p^2\to 0,\qquad \text{且}\quad \vec p\text{ 与 }\vec n\text{ 夹角很小}.
\end{equation}
共线胶子的发射在 QCD 中给出 $\int \dd\theta^2/\theta^2$ 型发散（在无质量理论中）。物理上，共线辐射被吸进部分子分布函数或喷射函数，不能用硬振幅单独描述。

\subsection{软}

软胶子四动量的所有分量都小：
\begin{equation}
  k^\mu\to 0.
\end{equation}
eikonal 近似下，发射软胶子的振幅因子化为
\begin{equation}
\label{eq:eikonal}
  \mathcal{M}(p;k)\approx g_s T^a\frac{p\cdot\varepsilon^a}{p\cdot k}\,\mathcal{M}(p).
\end{equation}
分母 $p\cdot k$ 在 $k$ 平行于 $p$ 时也小，所以软与共线相空间有重叠。这正是双对数的来源，也是 SCET 里 Wilson 线出现的原因：eikonal 因子是 Wilson 线的展开。

\subsection{硬}

硬模式的动量所有分量都 $\sim Q$，虚度 $p^2\sim Q^2$。它们不能在末态以实粒子出现（除非被积分进 Wilson 系数），在 EFT 里被积掉。

\begin{figure}[h]
\centering
\begin{tikzpicture}[scale=0.95]
  \draw[-{Stealth},thick] (-0.2,0) -- (6.4,0) node[right] {$k^-$};
  \draw[-{Stealth},thick] (0,-0.2) -- (0,4.2) node[above] {$k^+$};
  \draw[thick,scblue,fill=scblue!15] (0.15,3.1) rectangle (1.35,3.9);
  \node[font=\small] at (0.75,3.5) {$\bar n$-共线};
  \draw[thick,scblue,fill=scblue!15] (3.3,0.15) rectangle (5.5,0.85);
  \node[font=\small] at (4.4,0.5) {$n$-共线};
  \draw[thick,orange,fill=orange!20] (0.15,0.15) rectangle (1.15,1.0);
  \node[font=\small] at (0.65,0.55) {超软};
  \draw[thick,purple,fill=purple!12] (2.0,1.6) rectangle (3.4,2.7);
  \node[font=\small] at (2.7,2.15) {硬};
  \node[font=\small,align=left] at (8.1,2.2) {示意：\\光锥平面上的\\模式分区};
\end{tikzpicture}
\caption{$n$-共线模式大分量是 $k^-$，$+$ 分量很小；超软两个光锥分量都小；硬模式活在中间。真正的分区还依赖你用 SCET$_{\mathrm I}$ 还是 SCET$_{\mathrm{II}}$。}
\end{figure}

%======================================================================
\section{光锥坐标}
\label{sec:lc}
%======================================================================

\subsection{参考向量}

取两个类光向量
\begin{equation}
\label{eq:nnbar}
  n^2=0,\qquad \bar n^2=0,\qquad n\cdot\bar n=2.
\end{equation}
标准选择：粒子沿 $+z$ 运动时
\begin{equation}
  n^\mu=(1,0,0,1),\qquad \bar n^\mu=(1,0,0,-1).
\end{equation}
任意四矢量分解为
\begin{equation}
\label{eq:decomp}
  p^\mu=\frac{\bar n\cdot p}{2}\,n^\mu+\frac{n\cdot p}{2}\,\bar n^\mu+p_\perp^\mu
  \equiv p_-^\mu+p_+^\mu+p_\perp^\mu.
\end{equation}
记号
\begin{equation}
  p^+\equiv n\cdot p,\qquad p^-\equiv \bar n\cdot p,
\end{equation}
于是（度规 $(+,-,-,-)$）
\begin{equation}
\label{eq:p2}
  p^2=p^+ p^--p_\perp^2.
\end{equation}
常把分量写成
\begin{equation}
  p^\mu=(p^+,p^-,p_\perp).
\end{equation}

\subsection{哪个分量大？请算一次}

沿 $+z$ 的超相对论粒子：$p^\mu\approx(E,0,0,E)$，$E\sim Q$。则
\begin{equation}
  n\cdot p=E-E=0,\qquad \bar n\cdot p=E+E=2E\sim Q.
\end{equation}
所以 \textbf{$n$-共线粒子的大分量是 $p^-=\bar n\cdot p\sim Q$，小分量是 $p^+=n\cdot p$}。若喷射有小质量 $M_J$，则 $p^+\sim M_J^2/Q$。

\begin{warnbox}
文献里 $p^+$ 有时定义为 $\bar n\cdot p$。本讲义固定：$p^+=n\cdot p$。听课第一件事是看老师的 $n^\mu$ 指向哪、大分量叫 $+$ 还是 $-$。对上记号之前不要抄幂次。
\end{warnbox}

\subsection{Dirac 结构：共线投影}

对 $n$-共线无质量费米子，$\slashed{p}\,u=0$ 意味着旋量几乎满足
\begin{equation}
  \frac{\slashn\slashnb}{4}u=u.
\end{equation}
定义投影算符
\begin{equation}
\label{eq:proj}
  P_n=\frac{\slashn\slashnb}{4},\qquad P_{\bar n}=\frac{\slashnb\slashn}{4},\qquad
  P_n+P_{\bar n}=1,\quad P_n^2=P_n.
\end{equation}
共线夸克场是
\begin{equation}
  \xi_n=P_n\psi.
\end{equation}
另一半 $P_{\bar n}\psi$ 是重的（或称小分量），会被运动方程消掉，类似于 HQET 里的小分量。这是 SCET 拉氏量看起来像二维 Dirac 理论的原因。

%======================================================================
\section{幂次计数}
\label{sec:power}
%======================================================================

引入小参数 $\lambda\ll 1$。不同物理量按 $\lambda$ 的多少次方标度，称为 power counting。SCET 有两套常用方案。

\subsection{SCET$_{\mathrm I}$（超软--共线）}

适用于存在中间共线标度、且软标度是其平方再除以 $Q$ 的过程，例如：推力（thrust）、$B\to X_s\gamma$ 的光子能谱、喷射质量。
\begin{equation}
  \lambda^2\sim \frac{M_J^2}{Q^2}\sim \frac{p_c^2}{Q^2}.
\end{equation}
模式标度（写成 $(p^+,p^-,p_\perp)$，相对 $Q$）：
\begin{center}
\begin{tabular}{ll}
\toprule
模式 & 标度 \\
\midrule
硬 (hard) & $(1,1,1)Q$ \\
$n$-共线 (collinear) & $(\lambda^2,1,\lambda)Q$ \\
超软 (ultrasoft) & $(\lambda^2,\lambda^2,\lambda^2)Q$ \\
\bottomrule
\end{tabular}
\end{center}
共线虚度 $p_c^2\sim \lambda^2 Q^2$，超软虚度 $p_{us}^2\sim \lambda^4 Q^2$。共线与超软的 $p^+$ 同阶，因此它们之间的相互作用必须小心展开（multipole 展开）。

\subsection{SCET$_{\mathrm{II}}$（软--共线）}

适用于 TMD、排他衰变 $B\to D\pi$ 等：软与共线的横动量同阶。
\begin{center}
\begin{tabular}{ll}
\toprule
模式 & 标度 \\
\midrule
$n$-共线 & $(\lambda^2,1,\lambda)Q$ \\
软 (soft) & $(\lambda,\lambda,\lambda)Q$ \\
\bottomrule
\end{tabular}
\end{center}
此时软与共线虚度同为 $\lambda^2 Q^2$，但不能用同一个动能项描述，需要额外的快速度（rapidity）重整化。

\begin{keybox}
课前只需记住：\\
\textbf{共线} = 一个光锥分量大、一个小、横动量中等；\\
\textbf{超软} = 所有分量都 $\sim\lambda^2 Q$；\\
\textbf{软} = 所有分量都 $\sim\lambda Q$。\\
老师说 SCET$_{\mathrm I}$ 还是 ${}_{\mathrm{II}}$，就是在选第二行还是第三行。
\end{keybox}

\subsection{场的幂次}

场的标度由传播子（两点函数）定。粗略地，$n$-共线夸克 $\xi_n\sim\lambda$，共线胶子的横分量 $A_{n\perp}\sim\lambda Q$，超软夸克 $\sim\lambda^3$，超软胶子 $A_{us}\sim\lambda^2 Q$。不必背全表，但要理解：\textbf{拉氏量每一项都有确定的 $\lambda$ 次数}，leading-power 拉氏量只保留最低次项。

\begin{example}[共线动量如何满足 $p^2\sim\lambda^2 Q^2$]
$p\sim(\lambda^2,1,\lambda)Q$，则 $p^+p^-\sim\lambda^2 Q^2$，$p_\perp^2\sim\lambda^2 Q^2$，两者同阶，故 $p^2$ 可以是 $\lambda^2 Q^2$ 而不是更小。这保证共线粒子可以轻微离壳，正好描述喷射质量。
\end{example}

%======================================================================
\section{SCET 的基本结构}
\label{sec:structure}
%======================================================================

\subsection{把 QCD 场拆开}

形式上
\begin{equation}
  \psi=\sum_{n\in\text{方向}}\xi_n+q_{us}\quad\text{（或再加软夸克）},\qquad
  A^\mu=\sum_n A_n^\mu+A_{us}^\mu.
\end{equation}
每个共线方向 $n$ 有自己的一套场。两个对向喷射需要 $\xi_n$ 与 $\xi_{\bar n}$。硬场不出现：它们的效应进入 Wilson 系数。

\subsection{共线夸克拉氏量（直观推导）}

从 QCD 的 $\bar\psi i\slashed{D}\psi$ 出发，只保留 $n$-共线夸克，并作投影 $\xi_n=P_n\psi$。用运动方程消去小分量后，leading-power 结果具有结构
\begin{equation}
\label{eq:Lxi}
  \mathcal{L}_{\xi_n}=\bar\xi_n\Bigl(in\cdot D_n+i\slashed{D}_{n\perp}\frac{1}{i\bar n\cdot D_n}i\slashed{D}_{n\perp}\Bigr)\frac{\slashnb}{2}\xi_n.
\end{equation}
解读：
\begin{itemize}[leftmargin=1.4em]
  \item $n\cdot D_n$ 是大方向上的协变导数，量级 $\lambda^2 Q$；
  \item $\slashed{D}_{n\perp}$ 是横导数，量级 $\lambda Q$；
  \item $1/(i\bar n\cdot D_n)$ 来自消去小分量，量级 $1/Q$。
\end{itemize}
两项同为 $\lambda^2$，属于 leading power。这与 HQET 消去小分量得到 $iv\cdot D$ 是同一类计算。

\subsection{Wilson 线：规范不变的灵魂}

共线夸克单独携带共线规范变换，不构成规范不变量。必须用 Wilson 线把规范相位拖到无穷远。沿 $\bar n$ 方向的共线 Wilson 线
\begin{equation}
\label{eq:Wn}
  W_n(x)=\mathrm{P}\exp\Bigl(ig_s\int_{-\infty}^{0}\dd s\,\bar n\cdot A_n(x+s\bar n)\Bigr)
\end{equation}
把 $\xi_n$ 组成
\begin{equation}
  \chi_n=W_n^\dagger\xi_n,
\end{equation}
$\chi_n$ 在共线规范变换下中性（仍可带超软/软规范变换）。超软相互作用则由超软 Wilson 线
\begin{equation}
  Y_n(x)=\mathrm{P}\exp\Bigl(ig_s\int_{-\infty}^{0}\dd s\,n\cdot A_{us}(x+sn)\Bigr)
\end{equation}
生成。eikonal 发射公式~\eqref{eq:eikonal} 正是 $Y_n$ 的微扰展开。

\begin{keybox}
看到 $W_n$ 与 $Y_n$ 不要慌。它们只做一件事：\textbf{让带颜色的共线/软场重新变成规范不变量}。因子化定理里，软函数几乎总是 Wilson 线的关联函数。
\end{keybox}

\subsection{硬--共线流与匹配}

QCD 中的电磁流 $\bar\psi\gamma^\mu\psi$，在 SCET 里匹配成
\begin{equation}
\label{eq:current}
  \bar\psi\gamma^\mu\psi
  = \int\dd\omega\,\dd\omega'\,C(\omega,\omega',\mu)\,
  \bar\chi_{\bar n}(\omega')\gamma^\mu_\perp\chi_n(\omega)+\cdots.
\end{equation}
$C$ 是硬 Wilson 系数，阶 $\alpha_s$ 的匹配从 QCD 与 SCET 的单圈图相减得到。$\cdots$ 是更高幂次算符。\textbf{因子化的硬函数来自 $|C|^2$。}

\subsection{两种常用形式语言}

\begin{itemize}[leftmargin=1.4em]
  \item \textbf{标签形式}（label formalism, Bauer--Stewart 等）：把共线大动量标成标签 $p^-$，位置只描述剩余小动量。
  \item \textbf{位置空间形式}（Neubert、Becher 等常用）：不引入标签，用投影与 multipole 展开。
\end{itemize}
物理相同。Neubert 的课更常走位置空间。若幻灯片出现 $\mathcal{P}$ 或标签导数，那是第一套语言。

%======================================================================
\section{因子化定理的物理图像}
\label{sec:fact}
%======================================================================

\subsection{一句话}

在 leading power，许多截面（或衰减谱）写成
\begin{equation}
\label{eq:fact}
  \sigma=H\otimes J\otimes S
  \quad\text{（再卷积 PDF 等）}.
\end{equation}
\begin{itemize}[leftmargin=1.4em]
  \item $H$：硬函数，尺度 $Q$，无红外发散；
  \item $J$：喷射/共线函数，尺度 $M_J$ 或 $p_\perp$；
  \item $S$：软函数，尺度 $M_J^2/Q$ 或 $p_\perp^2/Q$。
\end{itemize}
每个函数只含一类模式，因而只含一种对数。卷积 $\otimes$ 来自共线与软之间共享的光锥动量分数或位置变量。

\subsection{为什么因子化``应该''成立}

SCET 拉氏量在 leading power 下，超软场只通过 $n\cdot A_{us}$ 进入，且可用场重定义把超软 Wilson 线从共线拉氏量里拆到流算符上。于是共线与超软在动力学上解耦：
\begin{equation}
  \mathcal{L}=\mathcal{L}_n+\mathcal{L}_{\bar n}+\mathcal{L}_{us}+\cdots,
\end{equation}
关联函数裂成硬系数乘共线矩阵元乘软矩阵元。这就是因子化定理的 EFT``证明提纲''。真正的证明还要处理快速度发散、 Glauber 模式、测量算符是否红外安全等；课前知道提纲即可。

\begin{example}[推力]
$e^+e^-$ 的推力 $T$ 接近 $1$ 时，$\tau=1-T\sim\lambda^2$。截面因子化为
\begin{equation}
  \frac{\dd\sigma}{\dd\tau}=H(Q,\mu)\int\dd s_n\,\dd s_{\bar n}\,\dd k\,
  J(s_n,\mu)J(s_{\bar n},\mu)S(k,\mu)\,\delta(\tau-(s_n+s_{\bar n}+k)/Q^2).
\end{equation}
三个函数分别吃硬、两个喷射、软辐射。这是 SCET$_{\mathrm I}$ 的展示橱窗。
\end{example}

\subsection{重整化群与 Sudakov 重求和}

每个函数满足 RG 方程。硬函数的典型形式含\textbf{cusp 反常}：
\begin{equation}
\label{eq:rge}
  \mu\frac{\dd}{\dd\mu}H(Q,\mu)=\Bigl[\Gamma_{\mathrm{cusp}}(\alpha_s)\log\frac{Q^2}{\mu^2}+\gamma_H(\alpha_s)\Bigr]H.
\end{equation}
$\Gamma_{\mathrm{cusp}}$ 来自 Wilson 线尖角的紫外发散，是双对数的源头。解这方程，等于把 $\alpha_s^n\log^{2n}$ 全部吃进指数（Sudakov 因子）。

操作步骤：
\begin{enumerate}[leftmargin=1.6em]
  \item 在各自的自然尺度算 $H,J,S$（那里对数小，固定阶可靠）；
  \item 用 RG 把它们演化到同一个 $\mu$；
  \item 卷积得到物理截面。
\end{enumerate}

\begin{hearbox}
Neubert 会强调：\textbf{重求和 = 解 RG}，不是另发明一套技巧。cusp 反常 $\Gamma_{\mathrm{cusp}}$ 是关键词。看到 $\log(Q^2/\mu^2)$ 乘在反常里，就知道双对数从哪里来。
\end{hearbox}

%======================================================================
\section{典型应用地图}
\label{sec:apps}
%======================================================================

五讲不可能覆盖全部应用，但下面几块最常出现。

\begin{enumerate}[leftmargin=1.8em]
  \item \textbf{重味}：$B\to X_s\gamma$、$B\to X_u\ell\nu$ 的端点谱。HQET 处理重夸克，SCET 处理衰变产物的共线与软胶子。
  \item \textbf{$e^+e^-$ 事件形状}：thrust、重喷射质量。精确 QCD 与 $\alpha_s$ 提取。
  \item \textbf{对撞机喷射}：喷射质量、$p_T$ 谱、窄喷射的因子化。与本暑期学校 Ian Moult 的能量关联器课程相邻。
  \item \textbf{TMD / SCET$_{\mathrm{II}}$}：横动量依赖 PDF，Drell--Yan 的 $q_T$ 谱。
  \item \textbf{次领头幂次}（power corrections）：leading-power 因子化之外的 $\lambda$ 修正，是近年热点。
\end{enumerate}

若课堂偏向重味，你会看到 $v^\mu$（重夸克速度）与 $n^\mu$ 同时出现；若偏向对撞机，会看到束流方向的共线 PDF。

%======================================================================
\section{课堂可能出现的公式与术语}
\label{sec:classroom}
%======================================================================

\begin{align}
  n^2&=\bar n^2=0,\quad n\cdot\bar n=2
  && \text{光锥基}\\
  p^\mu&=\tfrac{n\cdot p}{2}\bar n^\mu+\tfrac{\bar n\cdot p}{2}n^\mu+p_\perp^\mu
  && \text{分解}\\
  p_n&\sim Q(\lambda^2,1,\lambda)
  && \text{$n$-共线}\\
  p_{us}&\sim Q(\lambda^2,\lambda^2,\lambda^2)
  && \text{超软}\\
  \xi_n&=\frac{\slashn\slashnb}{4}\psi
  && \text{共线夸克}\\
  W_n&=\mathrm{P}\exp\bigl(ig\int \bar n\cdot A_n\bigr)
  && \text{共线 Wilson 线}\\
  \sigma&=H\otimes J\otimes S
  && \text{因子化}\\
  \mu\partial_\mu H&=\bigl[\Gamma_{\mathrm{cusp}}\log\tfrac{Q^2}{\mu^2}+\gamma\bigr]H
  && \text{硬函数 RG}
\end{align}

高频词：
\textit{scale separation, matching, Wilson coefficient, collinear / soft / ultrasoft, light-cone coordinates, power counting, leading power, Wilson line, eikonal, factorization, jet function, soft function, cusp anomalous dimension, Sudakov resummation, SCET$_{\mathrm I}$ / SCET$_{\mathrm{II}}$, rapidity divergence.}

%======================================================================
\section{常见困惑}
\label{sec:faq}
%======================================================================

\begin{enumerate}[leftmargin=1.8em]
  \item \textbf{$n$ 与 $\bar n$ 搞反，进而把 $p^\pm$ 搞反。}\\
  用一个具体四动量 $(E,0,0,E)$ 代入 $n\cdot p$ 与 $\bar n\cdot p$，立刻能检查。

  \item \textbf{软与超软当同义词。}\\
  在 SCET 里它们标度不同，对应 SCET$_{\mathrm{II}}$ 与 SCET$_{\mathrm I}$。口语里有人混用，看幂次表。

  \item \textbf{以为 SCET 是``另一种 QCD 微扰论''。}\\
  它是有效理论：硬模式被积掉，展开参数是 $\lambda$ 不是只靠 $\alpha_s$。$\alpha_s$ 仍在每个区里做圈图。

  \item \textbf{因子化 = 截面变成三个数相乘。}\\
  通常是卷积，不是普通乘法。自变量是质量、光锥分数或冲击参数。

  \item \textbf{Wilson 线是可选项。}\\
  没有它，共线场的规范不变性不完整，软辐射的 eikonal 结构也无法系统写出。

  \item \textbf{leading power 已经精确。}\\
  它是 $\lambda$ 的领头项。对数可以全部重求，但幂次修正（$M_J/Q$ 等）是另一层。

  \item \textbf{SCET 自动处理所有红外安全问题。}\\
  Glauber 胶子、非全局对数、非红外安全测量，都可能破坏最简单的因子化。课堂若提到，记下``假设条件''即可。

  \item \textbf{把 HQET 与 SCET 当成互斥。}\\
  重夸克用 HQET 场 $h_v$，轻衰变产物用 SCET 场。许多 $B$ 物理公式两者同时出现。
\end{enumerate}

%======================================================================
\section{听课重点}
\label{sec:listen}
%======================================================================

\begin{enumerate}[leftmargin=1.8em]
  \item \textbf{第一讲}：问题从哪来？Sudakov 双对数？事件形状？记下尺度层次~\eqref{eq:scales}。
  \item \textbf{第二讲}：光锥分解与 $\lambda$ 计数。对照本讲义第~\ref{sec:lc}--\ref{sec:power}~节，抄老师的表，不要抄本讲义的表（记号可能不同）。
  \item \textbf{第三讲}：拉氏量、投影、Wilson 线。问自己：没有 $W_n$ 哪一项不规范不变？
  \item \textbf{第四、五讲}：匹配、因子化、RG。把~\eqref{eq:fact} 与~\eqref{eq:rge} 对上具体过程。
  \item 讨论时段优先问：本课用 SCET$_{\mathrm I}$ 还是 ${}_{\mathrm{II}}$？次领头幂次是否讲？与能量关联器课程如何衔接？
\end{enumerate}

\begin{hearbox}
五讲的最小闭环：\\
\emph{尺度层次} $\to$ \emph{模式} $\to$ \emph{$\lambda$ 计数} $\to$ \emph{Wilson 线} $\to$ \emph{$H\otimes J\otimes S$} $\to$ \emph{cusp RG}。\\
能用自己的话把这六块串起来，预习就成功了。细节公式课堂上再填。
\end{hearbox}

%======================================================================
\section{进一步学习路线}
\label{sec:next}
%======================================================================

\subsection{课前（强烈建议做的两道小题）}

\begin{enumerate}[leftmargin=1.6em]
  \item 取 $p^\mu=(E,0,0,E-\Delta)$，$\Delta\ll E$，按~\eqref{eq:nnbar} 计算 $p^\pm,p_\perp,p^2$，并与 $(\lambda^2,1,\lambda)Q$ 对照，读出 $\lambda^2\sim\Delta/E$。
  \item 从 eikonal 顶点 $p\cdot\varepsilon/(p\cdot k)$ 看出它是 Wilson 线展开的第一项（只求感性认识）。
\end{enumerate}

\subsection{课堂同步}

\begin{itemize}[leftmargin=1.4em]
  \item Becher, Broggio, Ferroglia, \emph{Introduction to Soft-Collinear Effective Theory}（Springer Lecture Notes, arXiv:1410.1892）--- 最贴近 Neubert 风格的系统教材。
  \item Bauer, Stewart 的 MIT OCW 8.851 SCET 讲义（标签形式）。
  \item Becher, Les Houches lectures, arXiv:1803.04310。
\end{itemize}

\subsection{课后}

\begin{itemize}[leftmargin=1.4em]
  \item 原始文献：Bauer, Fleming, Pirjol, Stewart 系列；Neubert 与合作者关于 SCET 位置空间与 $B$ 物理的论文。
  \item 应用：thrust 因子化、喷射质量、TMD。
  \item 若听了 Moult 的能量关联器课：用 SCET 语言重读能量关联器的共线极限。
\end{itemize}

%======================================================================
\section{英文术语表}
\label{sec:glossary}
%======================================================================

\begin{center}
\begin{tabular}{>{\raggedright\arraybackslash}p{0.50\textwidth} p{0.42\textwidth}}
\toprule
\textbf{English} & \textbf{中文} \\
\midrule
soft-collinear effective theory (SCET) & 软共线有效理论 \\
effective field theory (EFT) & 有效场论 \\
scale separation & 尺度分离 \\
matching; Wilson coefficient & 匹配；威尔逊系数 \\
operator product expansion (OPE) & 算符乘积展开 \\
hard / collinear / soft / ultrasoft mode & 硬 / 共线 / 软 / 超软模式 \\
light-cone coordinates & 光锥坐标 \\
light-like vector & 类光矢量 \\
transverse momentum & 横动量 \\
power counting; leading power & 幂次计数；领头幂次 \\
Sudakov form factor / logarithm & 苏达科夫形状因子 / 对数 \\
resummation & 重求和 \\
eikonal approximation & 几何光学（eikonal）近似 \\
Wilson line; path ordering & 威尔逊线；路径排序 \\
collinear projection & 共线投影 \\
label formalism vs.\ position-space SCET & 标签形式 vs.\ 位置空间 SCET \\
factorization theorem & 因子化定理 \\
hard / jet / soft function & 硬 / 喷射 / 软函数 \\
cusp anomalous dimension & 尖角反常维数 \\
renormalization-group equation (RGE) & 重整化群方程 \\
rapidity divergence & 快速度发散 \\
SCET$_{\mathrm I}$ / SCET$_{\mathrm{II}}$ & 两套模式方案 \\
thrust; event shape & 推力；事件形状 \\
jet mass; PDF; TMD & 喷射质量；部分子分布；TMD \\
power correction & 幂次修正 \\
\bottomrule
\end{tabular}
\end{center}

\vspace{1.2em}
\noindent\textit{声明：}本讲义为课前预习材料，不是 Neubert 教授的正式讲义。光锥记号在文献中不统一，一律以课堂当场定义为准。

\end{document}
